\chapter{Methods}\label{methods}

\begin{figure}
{\centering
    \includegraphics[width=\textwidth]{study-area.pdf}
    \caption[Study Area: \acl{OM} Summit]{Study area at the summit of \acl{OM} in ArcGIS Pro. Sinusoidal Martian Projection. Contours in \unit{km}. Square is \qtyproduct{200 x 200}{\km}, centered at the midpoint of the \qty{19}{\km} contour.}
\label{study-area}
}
\end{figure}

\section{Overview}
Goal (a) of this project is to constrain the degree and character of 3D surface deformation recorded at the summit of \ac{OM} to the greatest spatial extent possible. See Figure~\ref{study-area} for the spatial extent of the study area. 3D deformation is a characterization of net change; thus, initial and final surface attitude conditions are necessary. Attitude is completely defined by two angles; strike and dip are most widely used in geology. In this inquiry, dip direction will be used instead of strike because the lava flows naturally point in the dip direction of the surface upon which they are emplaced. No assumptions are made \textit{a priori} about large-scale spatial patterns in 3D deformation; therefore, independent calculations will be made at all locations where sufficient information is available.

Goal (b) of the project is to estimate magma reservoir inflation conditions (location, depth, size of reservoir, degree of inflation) consistent with observed 3D deformation at the surface. This process will entail running a series of numerical models simulating various possible combinations of magma reservoir conditions, analyzing the resultant surface deformation associated with each, and identifying the best match with the observed and calculated surface parameters.

While the technical definitions of these surface parameters will be presented later in Section~\ref{paleosurface}, it is worth noting here that 3D deformation discussed previously is a scalar term: in particular, the angular distance between the two normal vectors corresponding to each of the initial and final surfaces. This parameter plays two key roles in this project.

First, relevant to goal (a), it builds on the qualitative observations of \textcite{mouginis-mark_late-stage_2019} and extends the quantitative horizontal description of flow misalignment in \textcite{chadwick_late_2015} by accounting for vertical deformation immediately as a central component of the map-derived data collection.\footnote{Vertical deformation is accounted for in the geophysical modeling stage of \textcite{chadwick_late_2015}.} In other words, constraints on 3D deformation may be computed directly from the \ac{DEM} and mapped lava flows \& channels, prior to running numerical models. This strategy gives first-order insights by mapping regions where especially high magnitudes of 3D deformation are required to explain observed surface features. 

However, the 3D deformation parameter cannot simply be calculated for each point because insufficient attitude information is available for the paleosurface upon which lava flows and channels were originally emplaced. Recall that 3D deformation is the angular distance between the two normal vectors defining each of the initial and final surfaces. Each normal vector (equivalently, each surface attitude) requires two variables: dip and dip direction. Both of these are available for the modern (final) surface \ac{DEM}; however, only dip direction of the paleosurface can be inferred directly from the direction of lava flows. Rather than attempting to estimate paleoslope (dip) from lava flow morphology, this parameter is \textit{constrained.} Specifically, the paleo-dip assigned to each point answers the question: \textit{given modern dip, modern dip direction, and paleo-dip direction, what paleo-dip minimizes 3D deformation?} Through the simplifying assumption of minimum 3D deformation, two variable constraints (paleo-dip \& 3D deformation) can be obtained for the price of one, so to speak.

Second, relevant to goal (b), the previous discussion has shown how an assumption regarding of 3D deformation yields an estimate of paleo-dip which is independent of lava flow morphology (and can thus be compared with estimates derived from such methods). This is important because the numerical models used to test magma reservoir inflation events require construction of a paleo-edifice which no longer exists, and whose geometry must thus be inferred. Therefore, estimates of paleo-dip in addition to measurements of paleo-dip direction will guide first attempts at constructing this paleo-edifice.

To summarize, the 3D deformation strategy is to \begin{inparaenum}[a)]
    \item identify suitable point locations across the summit,
    \item collect attitude information for the modern surface at each point,
    \item infer paleoattitude information based on lava flows or channels at each point, and
    \item constrain total 3D deformation based on initial and final surface attitudes at each point. In this final step, an estimate of paleo-dip is produced as a additional useful result.
\end{inparaenum}
The following sections describe the details of each of these steps.

\section{Preparation of Published Data}
The \ac{CTX}\footnote{aboard the \ac{MRO} spacecraft launched in 2005 by \ac{NASA}.} captures \qty{\sim30}{\km} swaths across the entire martian surface in visible greyscale at \qty{\sim6}{\m} resolution. \textcite{Dickson2018AGB} blended these swaths to produce a raster mosaic product (hereafter, ``\ac{CTX} mosaic'') which is used here to visually identify lava flows and flow channels. In other words, this is the source from which paleo-dip direction is estimated across the summit region.

\ac{MOLA}\footnote{aboard the now-retired \ac{MGS} spacecraft launched in 1996 by \ac{NASA}.} returned topography data with horizontal resolution of \qtyproduct{300 x 1000}{\m} at the equator (better at high latitudes) and elevation uncertainty of \qty{\sim3}{\m}~\parencite{smith_mars_2001}. To improve spatial resolution, additional elevation data from the \ac{HRSC}\footnote{aboard the \ac{MEX} spacecraft launched in 2003 by the \ac{ESA}} was blended to product a \ac{DEM} with \qty{200}{\m} pixel resolution. Each pixel's vertical uncertainty is \qty{\sim1}{\m}, with an additional global uncertainty of \qty{\sim1.8}{\m} in the martian areoid (martian equivalent of Earth's geoid). In this project, the global areoid uncertainty is not a concern because only one region (the summit of \ac{OM}) is considered.

These two data sources were registered in an equal-area sinusoidal Mars projection in ArcGIS Pro. The study area is defined by a square \qtyproduct{200 x 200}{\km} centered at the centroid of the \qty{19}{\km} contour.\footnote{This is the highest integer \unit{km} which is roughly circular and completely encloses the caldera complex, implying that it largely records the conical shape of the shield edifice without influence from subsequent caldera collapse or reservoir inflation.} See Figure~\ref{study-area}.

\section{Mapping \& Sampling Site Selection}\label{mapping}
The \ac{CTX} mosaic was used to identify lava flows near the summit of \ac{OM}. Where possible, lobate flows were outlined and mapped as polygons. Flow centerlines were derived from these polygons. See Figure~\ref{flow}. Channels were mapped as lines where flow margins were not visible. Some channels were discontinuous; these features were inferred to be regions where channels formed lava tubes which have subsequently collapsed to form a chain of skylights. See Figure~\ref{channel}.

\begin{figure}
{\centering
    \includegraphics[width=\textwidth]{flow.pdf}
    \caption[Mapped lava flow \& centerline]{A lobate flow is identified by relatively clear boundaries at each margin and mapped as a polygon (white) from which a centerline is derived for sampling.}
\label{flow}
}
\end{figure}

\begin{figure}
{\centering
    \includegraphics[width=\textwidth]{channel.pdf}
    \caption[Mapped lava channel]{A lava channel is mapped as a linear feature, including regions of discontinuity which are inferred to be collapsed skylight chains over lava tubes.}
\label{channel}
}
\end{figure}

\begin{figure}
{\centering
        \includegraphics[width=\textwidth]{sampling.pdf}
        \caption[Sampling site selection]{Each linear feature is assigned an average paleo-dip direction. Points are selected for sampling and calculations at \qty{5}{\km} and \qty{3}{\km} intervals for flows and channels, respectively. Paleo-dip direction is assigned to each point from its corresponding line; modern dip and dip direction is assigned to each point from its unique \ac{DEM} neighborhood.}
\label{sampling}
}
\end{figure}

In both cases, lava features are ultimately mapped as lines\footnote{Lines can have more than one vertex, i.e., they can be curved.} (rather than polygons) in ArcGIS Pro. While the \ac{DEM} covers the entire summit region, it is only along these lines that information about the paleosurface can be inferred. Therefore, points along these lines are selected as sampling locations where attitude measurements and calculations will take place. Points are produced at evenly spaced intervals of \qty{5}{\km} along each flow centerline and \qty{3}{\km} along each channel. No points are assigned to linear features of less than this length to ensure that flows and channels are long enough to record large-scale topography rather than local variation.

\section{Constraining 3D Deformation}
\subsection{Surface Attitude Data Collection}
The ``Surface Parameters'' tool was used in ArcGIS Pro to compute average \acf{sl} and \acf{az} across the entire study area. In both cases, a new raster is produced at the same spatial resolution and coverage as the original \ac{DEM}. To avoid capturing local topographic anomalies, these values are averaged\footnote{Technically, a quadratic surface is interpolated over the neighborhood region and dip and dip direction are computed at the center of this surface.} over a circular ``neighborhood'' with radius \qty{2}{\km}. To each sampling location identified in Section~\ref{mapping} is assigned \acf{sl} and \acf{az} based on the value of the corresponding raster at that location.

\subsection{Constraining Paleo-slope}
Consider a point on the surface of a unit sphere. It is uniquely described by three spherical coordinates $(\acs{spherical-radial},\acs{spherical-azimuth},\acs{spherical-zenith})$. Since the sphere is of unit radius, $\acs{spherical-radial}=1$. Therefore, it is sufficient to define a point $\acs{point}\equiv(\acs{spherical-azimuth}_p,\acs{spherical-zenith}_p)$.

Now consider a line of constant \acs{spherical-azimuth} on the unit sphere. It passes through the zenith (pole) of the sphere and is analogous to a line of longitude on the Earth. This line can be defined $\acs{line}\equiv(\acs{spherical-azimuth}_\ell)$. \emph{Which point \acs{closest-point} on \acs{line} is closest to \acs{point}?}

\begin{figure}
    % view angle ({zenith}{azimuth})
\pgfmathsetmacro{\viewze}{65}
\pgfmathsetmacro{\viewaz}{105}
\tdplotsetmaincoords{\viewze}{\viewaz}

% define lengths
\pgfmathsetmacro{\radius}{1.3}
\pgfmathsetmacro{\axislength}{1.4}

% define angles for p
\pgfmathsetmacro{\az}{40}
\pgfmathsetmacro{\ze}{37}

% define angles for p' (no ' in macros!)
\pgfmathsetmacro{\azi}{330}
\pgfmathsetmacro{\zen}{14.4523587}

\pgfmathsetmacro{\cenang}{34.438487775}

\begin{tikzpicture}[scale=5,tdplot_main_coords]

% origin
\coordinate (O) at (0,0,0);

% also defines (Pxy), (Pxz), (Pyz), etc.
\tdplotsetcoord{P}{\radius}{\ze}{\az}
\tdplotsetcoord{P'}{\radius}{\zen}{\azi}
\tdplotsetcoord{Q}{\radius}{90}{\azi}

% grey hemisphere
\filldraw[tdplot_screen_coords, color = gray, opacity = 0.1] (0,0) -- (\radius,0) arc (0:180:\radius) -- (0,0);
\filldraw[tdplot_main_coords, color = gray!10!white] (0,0) -- (0:\radius) arc (0:365:\radius) -- (0,0);

% circumference
\tdplotdrawarc[very thin, dashed]{(O)}{\radius}{\viewaz}{\viewaz+180}{}{}
\tdplotdrawarc[very thin]{(O)}{\radius}{\viewaz-180}{\viewaz}{}{}

% az' surface
\tdplotsetthetaplanecoords{\azi}

% fill az' surface
\filldraw[tdplot_rotated_coords, color = red, opacity = 0.45] (\radius,0) arc (0:\zen:\radius) -- (0,0);
\filldraw[tdplot_rotated_coords, color = black, opacity = 0.05] (\radius,0) arc (0:90:\radius) -- (0,0);

% line az' surface
\draw[very thin] (O) -- (P') node[anchor = south east] {\acs{closest-point}};
\draw[very thin] (O) -- (Q);
\tdplotdrawarc[tdplot_rotated_coords, very thick]{(O)}{\radius}{90}{14}{anchor=north west}{\acs{line}}
\tdplotdrawarc[tdplot_rotated_coords,very thick,red]{(O)}{\radius}{0}{\zen}{anchor=south}{\acs{ze'}}

% az surface
\tdplotsetthetaplanecoords{\az}

% line az surface
\draw[very thin] (P) -- (O) -- (Pxy) -- (P) node[anchor = south west] {\acs{point}};

% cartesian axes
\draw[arrow] (O) -- (\axislength,0,0) node[anchor=north]{$x$};
\draw[arrow] (O) -- (0,\axislength,0) node[anchor=west]{$y$};
\draw[arrow] (O) -- (0,0,\axislength) node[anchor=south]{$z$};

% az' angle label
\tdplotdrawarc{(O)}{0.15}{0}{\azi}{anchor=south west}{\acs{az'}}

% az angle label
\tdplotdrawarc{(O)}{0.25}{0}{\az}{anchor=north}{\acs{az}}

% ze angle label
\tdplotdrawarc[tdplot_rotated_coords]{(O)}{0.3}{0}{\ze}{anchor=south west}{\acs{ze}}

% central angle plane
\tdplotsetrotatedcoords{\azi}{-90+\zen}{0}

% central angle line
\tdplotdrawarc[tdplot_rotated_coords,very thick,blue]{(0,0,0)}{\radius}{\cenang}{0}{anchor=south west}{\acs{central-angle}}

% central angle fill
\filldraw[tdplot_rotated_coords, color = blue, opacity = 0.15] (\radius,0) arc (0:\cenang:\radius) -- (0,0);

% draw points p and p' 
\fill (P) circle [radius=0.4pt];
\fill[red] (P') circle [radius=0.4pt];

\end{tikzpicture}
    \caption[Constraint]{Constraint}
    \label{constraint}
\end{figure}

\Acf{paz} is assigned to each point directly from the linear lava feature from which the point was sampled. Therefore, it is important to ensure that each feature is oriented with the correct sense, i.e., pointing downhill rather than uphill. Longer features with noticeable large-scale bends are broken into multiple pieces which are approximately straight. Then, the primary dip direction is calculated for each feature and assigned to all points sampled within it.

The \ac{central-angle} subtended between two points on a unit sphere at the sphere's center is:
\begin{equation}
    \ac{central-angle}=\arccos\left[\cos\ac{spherical-zenith}_1\cos\ac{spherical-zenith}_2+\sin\ac{spherical-zenith}_1\sin\ac{spherical-zenith}_2\cos\left(\acs{spherical-azimuth}_2-\acs{spherical-azimuth}_1\right)\right].\label{gcd}
\end{equation}
Given $\ac{spherical-azimuth}_p$, $\ac{spherical-zenith}_p$, and $\ac{spherical-azimuth}_q$, which value of $\ac{spherical-zenith}_q$ minimizes \ac{central-angle}? This optimization problem can be solved by setting:
\begin{equation}
    \frac{\partial}{\partial \ac{spherical-zenith}_q}\arccos\left[\cos\ac{spherical-zenith}_q\cos\ac{spherical-zenith}_p+\sin\ac{spherical-zenith}_q\sin\ac{spherical-zenith}_p\cos\left(\acs{spherical-azimuth}_p-\acs{spherical-azimuth}_q\right)\right]=0.
\end{equation}
Collecting constant terms allows for the following simplification:
\begin{equation}
    \frac{\partial}{\partial \ac{spherical-zenith}_q}\arccos\left[a\cos\ac{spherical-zenith}_q+b\sin\ac{spherical-zenith}_q\right]=0,
\end{equation}
where $a=\cos\ac{spherical-zenith}_p$ and $b=\sin\ac{spherical-zenith}_p\cos\left(\acs{spherical-azimuth}_p-\acs{spherical-azimuth}_q\right)$. Differentiating using the chain rule gives:
\begin{equation}
    \frac{-1}{\sqrt{1-{\left(a\cos\ac{spherical-zenith}_q+b\sin\ac{spherical-zenith}_q\right)}^2}}\cdot\left(-a\sin\ac{spherical-zenith}_q+b\cos\ac{spherical-zenith}_q\right)=0.
\end{equation}
Clearly, the first term cannot be equal to $0$. Therefore:
\begin{gather}
    -a\sin\ac{spherical-zenith}_q+b\cos\ac{spherical-zenith}_q=0,\nonumber\\
    b\cos\ac{spherical-zenith}_q=a\sin\ac{spherical-zenith}_q,\nonumber\\
    b/a=\tan\ac{spherical-zenith}_q,\nonumber\\
    \ac{spherical-zenith}_q=\arctan(b/a).
\end{gather}
Substituting the constant terms back in:
\begin{gather}
    \ac{spherical-zenith}_q=\arctan\left[\frac{\sin\ac{spherical-zenith}_p\cos\left(\acs{spherical-azimuth}_p-\acs{spherical-azimuth}_q\right)}{\cos\ac{spherical-zenith}_p}\right],\nonumber\\
    \ac{spherical-zenith}_q=\arctan\left[\tan\ac{spherical-zenith}_p\cos\left(\acs{spherical-azimuth}_p-\acs{spherical-azimuth}_q\right)\right].
\end{gather}

\section{Crater Dating Lava Flows}

\section{Axisymmetric Model of Reservoir Inflation}
The \ac{n} is constructed for each surface in cartesian coordinates as follows:
\begin{equation}
    \acs{n} = \begin{bmatrix}
        \sin\acs{sl}\cos\acs{az}\\
        \sin\acs{sl}\sin\acs{az}\\
        \cos\acs{sl}
    \end{bmatrix}\qquad
    \acs{n}' = \begin{bmatrix}
        \sin\acs{psl}\cos\acs{paz}\\
        \sin\acs{psl}\sin\acs{paz}\\
        \cos\acs{psl}
    \end{bmatrix}
\end{equation}
The \ac{d} is used to describe the minimal 3D rotation necessary to deform the paleosurface attitde into the modern surface observed today. Mathematically, a \ac{R} is desired to satisfy
\begin{equation}
    \acs{R}\acs{n}'=\acs{n}.
\end{equation}
The explicit terms of \acs{R} are derived from the axis-angle form of the 3D rotation.
\begin{equation}
    \acs{d} = \acs{n}' \times \acs{n}.
\end{equation}

\subsection{Construction of Proto-Edifice}
\subsection{Independent Variables Tested}
